%============================ OBJETIVOS (ENTREGA 1) =========================
\chapter{Objetivos}
\label{cap:objetivos}

\section{Objetivo general}
Diseñar, implementar y validar un sistema de bajo costo, denominado EILOR, para el monitoreo en tiempo real del espectro de 2.4~GHz, la identificación de dispositivos y amenazas inalámbricas, y la gestión responsable del acceso de invitados mediante un punto de acceso con portal cautivo de registro, integrando hardware ESP32 y un servidor web de tiempo real.

\section{Objetivos específicos}
\begin{enumerate}[label=\textbf{OE\arabic*.},leftmargin=2.4em]
  \item Desarrollar un firmware para ESP32 que escanee redes Wi\nobreakdash-Fi (canales 1--13) y balizas BLE (canales 37--39) en modos dedicados conmutables y transmita los resultados en tiempo real mediante WebSocket.
  \item Construir un servidor modular en Node.js que agregue los datos, calcule un puntaje de interferencia por canal, recomiende el canal óptimo y detecte puntos de acceso de tipo gemelo malicioso.
  \item Implementar la identificación de fabricante por OUI y la detección de direcciones MAC aleatorizadas para enriquecer el inventario de dispositivos.
  \item Incorporar una capa de seguridad con autenticación por token de sesión para el control y un token opcional de dispositivo para la ingesta de datos.
  \item Desarrollar un punto de acceso (SoftAP) con portal cautivo que registre invitados (nombre, apellido y teléfono) y persista la información en el servidor, con reenvío ante caídas de enlace.
  \item Validar funcionalmente el sistema completo mediante pruebas de ingesta, análisis, persistencia, control de acceso y registro de invitados.
\end{enumerate}

\section{Matriz de objetivos bajo criterios SMART}
Cada objetivo específico se formuló siguiendo los criterios SMART (específico, medible, alcanzable, relevante y temporal). La Tabla~\ref{tab:smart} presenta la matriz de verificación.

\begin{table}[H]
\centering
\caption{Matriz SMART de los objetivos específicos}
\label{tab:smart}
\footnotesize
\begin{tabularx}{\linewidth}{@{}c X X l@{}}
\toprule
\textbf{OE} & \textbf{Indicador medible} & \textbf{Criterio de logro (alcanzable/relevante)} & \textbf{Plazo} \\
\midrule
OE1 & N.º de canales escaneados y latencia de ciclo & 13 canales Wi-Fi + 3 BLE; ciclo $\le$ 3~s & Sem. 6 \\
OE2 & Puntaje por canal y canal recomendado & Cálculo para 13 canales y salida 1/6/11 & Sem. 7 \\
OE3 & Cobertura del catálogo OUI & $\ge$ 120 prefijos; detección de MAC aleatoria & Sem. 7 \\
OE4 & Respuestas de control autenticado & 401 sin token / 200 con token & Sem. 7 \\
OE5 & Registros entregados por el portal & 100\,\% de envíos con enlace; \emph{buffer} sin enlace & Sem. 7 \\
OE6 & Casos de prueba superados & $\ge$ 90\,\% de casos funcionales OK & Sem. 8 \\
\bottomrule
\end{tabularx}
\end{table}

\noindent La justificación de cada objetivo se articula con las preguntas específicas del Capítulo~\ref{cap:problema}: OE1 y OE2 responden a la necesidad de diagnóstico del espectro; OE3 y OE6 a la identificación de dispositivos y su validación; OE4 y OE5 a la gestión responsable y segura del acceso de invitados.
